\documentclass[11pt]{article}
\usepackage[utf8]{inputenc}
\usepackage[T1]{fontenc}
\usepackage{geometry}
\geometry{a4paper, margin=1in}
\usepackage{graphicx}
\usepackage{amsmath, amssymb}
\usepackage{lineno}
\usepackage{setspace}
\usepackage[superscript,biblabel]{cite}

\title{A computational framework predicting adolescent scoliosis onset}

\author{Sayuj Krishnan S., MBBS, DNB (Neurosurgery) \\
\textit{Yashoda Hospitals, Hyderabad, India} \\
Correspondence: hellodr@drsayuj.info}

\date{}

\begin{document}

\maketitle

\textbf{Study Design.} Computational modeling and cross-species scaling analysis validated against historical clinical cohorts.

\textbf{Objective.} To develop a predictive framework for Adolescent Idiopathic Scoliosis (AIS) onset based on the mismatch between spinal growth velocity and metabolic supply scaling.

\textbf{Summary of Background Data.} Current theories of AIS focus on mechanical instability or genetic factors in isolation. However, the unique susceptibility of humans to AIS during the adolescent growth spurt suggests a system-level failure where rapid skeletal growth outpaces the scaling capacity of neuromuscular maintenance.

\textbf{Methods.} We developed a 3D Cosserat rod model of the spine coupled to a "Developmental Information Field" representing neuromuscular control. We derived an "Energy Deficit Window" based on allometric scaling laws, where mechanical demand scales as $L^4$ while metabolic supply scales as $L^2$. This model was validated against cross-species data (n=9) and compared with clinical progression rates from Weinstein (1983) and Lonstein (1984). Structural anisotropy of 23 paraspinal proteins was analyzed using AlphaFold to identify molecular failure points.

\textbf{Results.} Our analysis identifies a critical spinal length ($L_{crit} \approx 0.35$ m) where metabolic demand exceeds supply, corresponding to the onset of the adolescent growth spurt. The model predicts a "Rescue Cliff" where mechanosensors with high structural anisotropy (e.g., Vimentin, Anisotropy > 7.0) fail to maintain calibration under energy deficit. This failure leads to a loss of proprioceptive fidelity, allowing small asymmetries to amplify into Cobb angles $>15^\circ$. The model accurately reproduces the 10:1 female prevalence of severe curves based on earlier peak height velocity and lower basal mitochondrial density in paraspinal muscles.

\textbf{Conclusions.} AIS can be understood as a "Metabolic Buckling" event driven by an energy deficit during rapid growth. This computational framework provides a quantitative explanation for the timing of onset and suggests that therapeutic interventions should target metabolic support for paraspinal muscles during the critical growth window.

\vspace{1em}

\section*{Introduction}

Adolescent Idiopathic Scoliosis (AIS) remains a complex deformity with a multifactorial etiology\cite{cheng2015}. While mechanical factors like the "vicious cycle" of asymmetric loading are well-established, the primary trigger for the initial loss of stability remains elusive. We propose that the unique vulnerability of the human spine during adolescence arises from a fundamental mismatch between the scaling laws of mechanical demand and metabolic supply.

In this study, we present a computational framework that integrates biomechanics with metabolic scaling theory. We define the "Energy Deficit Window," a developmental period where the energetic cost of maintaining spinal alignment exceeds the available metabolic power. This deficit compromises the function of high-maintenance mechanosensory proteins, leading to a loss of proprioceptive control and subsequent buckling.

\section*{Methods}

\subsection*{Computational Model}
We utilized \texttt{PyElastica}\cite{tekinalp2023}, an open-source Python implementation of Cosserat rod theory, to model the spine as a continuous elastic rod. The rod was discretized into $N=50-100$ elements. The control system was implemented via an Information-Elasticity Coupling (IEC) term that modifies the local rest curvature $\bm{\kappa}^0(s)$ based on a scalar information field $I(s)$ representing HOX patterning.

\subsection*{Energy Deficit Calculation}
The metabolic cost of maintaining alignment ($P_{counter}$) was defined as proportional to $L^2 \int (\kappa_{IEC} - \kappa_{passive})^2 ds$. Supply capacity ($S_{proprio}$) was modeled scaling as $L^2$ (surface-limited) or $L^3$ (volume-limited). The deficit ratio $R = P_{counter} / S_{proprio}$ was computed across spinal lengths $L \in [0.2, 0.6]$ m to identify the critical crossover point.

\subsection*{Clinical Validation}
We validated the model predictions against historical cohort data. The predicted onset of instability was compared with curve progression rates reported by Weinstein\cite{weinstein1983} and Lonstein\cite{lonstein1984}. We also analyzed the sexual dimorphism in curve progression by incorporating sex-specific growth velocities and muscle mass ratios.

\section*{Results}

\subsection*{The Energy Deficit Window}
Our scaling analysis reveals that humans occupy a unique "Allometric Trap" ($Bg < 0.1$) where passive stiffness is insufficient to maintain alignment against gravity. The calculated energy deficit window begins at a spinal length of $L \approx 0.35$ m, coinciding with the early adolescent growth spurt. At $L=0.45$ m, the deficit reaches peak intensity, correlating with the highest risk of curve progression.

\subsection*{Protein Structural Failure}
Analysis of 23 key paraspinal proteins revealed a "Demand-Supply Anisotropy Gap" of 34\%. Mechanosensors required for alignment (e.g., Vimentin, PIEZO2) are structurally expensive (high anisotropy), while metabolic regulators (PPARGC1A) are intrinsically disordered. This structural mismatch suggests that under energy stress, the sensory system fails before the metabolic supply can adapt.

\section*{Discussion}

Our results support the hypothesis that AIS is a system failure driven by a transient energy deficit. The "Metabolic Buckling" framework explains why scoliosis is specific to the rapid growth phase in humans and offers a new perspective on the sexual dimorphism of the disease.

\begin{thebibliography}{10}
\bibitem{cheng2015} Cheng, J.C. et al. Adolescent idiopathic scoliosis. \textit{Nat. Rev. Dis. Primers} \textbf{1}, 15030 (2015).
\bibitem{tekinalp2023} Tekinalp, A. et al. PyElastica: Open-source software for the simulation of assemblies of slender, one-dimensional structures using Cosserat rod theory. \textit{Zenodo} (2023).
\bibitem{weinstein1983} Weinstein, S.L. & Ponseti, I.V. Curve progression in idiopathic scoliosis. \textit{J. Bone Joint Surg. Am.} \textbf{65}, 447-455 (1983).
\bibitem{lonstein1984} Lonstein, J.E. & Carlson, J.M. The prediction of curve progression in untreated idiopathic scoliosis during growth. \textit{J. Bone Joint Surg. Am.} \textbf{66}, 1061-1071 (1984).
\end{thebibliography}

\end{document}
